\section{The Same Person, Raised in Glory}

The soul's immediate blessedness does not exhaust the human destiny. God created us as embodied persons, and the salvation accomplished by Christ reaches the body. The resurrection of the dead belongs to the Creed because the body belongs to the person whom God saves. Our final hope is the restoration and glorification of the whole human being. The separated soul truly lives with Christ; nevertheless, the common consummation still awaits his return, the raising of the dead, and the manifestation of his victory throughout creation. Heaven's present reality and the future resurrection are parts of one promise.

Paul grounds this promise in an event: Christ has risen from the dead, the first fruits of those who have fallen asleep (1~Cor 15:20--28). The harvest takes its character from him. He has entered the condition toward which his members are being led, and death itself will be destroyed. The Christian therefore looks to the risen Lord when considering his own body. The body's future is not determined by the weakness visible in it now. Its final destiny rests upon the power of Christ, who will conform our lowly body to his glorious body (Phil 3:20--21).

Paul compares the body to a seed sown in the earth (1~Cor 15:35--44). The seed dies, and what grows differs greatly from what was sown. God gives it a new manner of life, and the contrast is immense: corruption yields to incorruption, dishonor to glory, weakness to power. The risen body is our own body, transformed by God. Its identity does not require the continuation of every present limitation. Resurrection means that the Creator is faithful to his creature while bringing that creature into a condition beyond its present powers. The likeness to a seed illustrates the change without explaining how God will raise the dead.

A ``spiritual body'' remains a body. Paul is naming the life-giving rule of the Spirit, not the replacement of human flesh by a disembodied existence. The bodily humanity of the risen Christ is the pattern of our hope (CCC 999--1000). In our present condition the body is vulnerable to exhaustion, disease, injury, and death; its desires can also resist the good known by reason. In glory bodily life will participate in the peace of the whole person. The body will express the soul's communion with God without the corruption and division that now burden it.

Irenaeus presses the reality of this hope through the Eucharist. In \emph{Against Heresies} V.2.2--3, he answers those who deny that flesh can receive salvation. The Lord shed real blood for us, and he gives his body and blood as food. The flesh nourished by this gift belongs to him and cannot be dismissed as intrinsically incapable of eternal life. Irenaeus joins the Creator's gifts of bread and wine to the Redeemer's gift of himself. The same divine generosity sustains our present bodies, gives the sacrament, and will grant resurrection to the bodies committed to the earth.

Irenaeus emphasizes that the body's incorruptibility comes from God, whose strength brings life out of our weakness. Immortality is received; the creature does not become the independent source of its own duration. The body's return to dust reveals how completely we depend upon the Creator. Resurrection displays what that Creator can give. Our dependence will endure in glory, acknowledged with gratitude by creatures who live through his generosity.

The traditional language of the glorified body's freedom from suffering and corruption, its glory, and its complete obedience to the life of the Spirit gathers the substance of this promise. We can consider this promise while caring for someone whose body now appears frail or unattractive. We cannot estimate his heavenly dignity by his present strength, intelligence, or independence. A body that requires washing, feeding, or patient assistance belongs to someone called to resurrection. Care for that body is therefore an act of reverence toward God's handiwork. The promise of glory deepens the seriousness of bodily mercy while illness remains.

This also changes how we regard our own vulnerability. Age or disability can take away a form of usefulness around which we built our identity. The loss is real, but the call to communion with God does not depend upon preserving that usefulness. We remain persons capable of receiving and giving love, even when the manner of that exchange becomes hidden from others. The final restoration is Christ's work. Our present bodily life, whether vigorous or greatly limited, can be entrusted to him and offered within the communion of his members. The body he will raise is the body he already claims as his own.

Paul concludes his teaching on resurrection by urging steadfast labor in the Lord, because such labor is not vain (1~Cor 15:58). That exhortation follows his account of the body's victory over death. The bodily needs of the people with whom we now live belong within that service. We are free to seek healing, relieve pain, honor the dead, and patiently bear what cannot yet be remedied. Beneath these differing actions stands one hope: Christ's victory will reach the whole person, and death will not retain what God has made for life.

